\section{РЕЗУЛЬТАТЫ ЭКСПЕРИМЕНТОВ И АНАЛИЗ}

В разделе приведены результаты двух серий экспериментов --- пакетной аналитики на
наборе TPC-DS и near-real-time-сценария при потоковой записи, --- их анализ по типам
запросов и рекомендации по применению форматов хранения.

\subsection{Сценарий 1: пакетная аналитика на наборе TPC-DS}

На наборе TPC-DS~\cite{tpcds-spec} масштаба \SI{100}{\giga\byte} сравнивались форматы Parquet и Vortex
по времени выполнения \num{103} запросов (таблицы только для добавления, крупные
файлы разделов, параллелизм 300). Из \num{103} запросов Vortex быстрее Parquet более
чем на \SI{10}{\percent} на \num{82} запросах; среднее ускорение по запросу ---
порядка \SI{25}{\percent}; пиковые ускорения достигают \SI{80}{\percent}. Parquet
быстрее Vortex лишь на пяти запросах. Характерные запросы приведены в
таблице~\ref{tab:tpcds}.

\begin{table}
    \caption{Характерные результаты на наборе TPC-DS (масштаб \SI{100}{\giga\byte}): относительное изменение времени выполнения Vortex по сравнению с Parquet}
    \label{tab:tpcds}
    \begin{tabularx}{\textwidth}{|l|c|X|}\hline
    \textbf{Запрос} & \textbf{Vortex к Parquet} & \textbf{Особенности запроса (объяснение)} \\ \hline
    q51 & $\approx -80\,\%$ & оконные нарастающие суммы по продажам, мало столбцов, селективные условия --- column pruning и нативное проталкивание предикатов работают в полную силу \\ \hline
    q42 & $\approx -52\,\%$ & агрегация продаж по категории с фильтром по дате --- селективный предикат + узкая проекция \\ \hline
    q35 & $\approx -52\,\%$ & демографический срез с подзапросами существования --- многократное сканирование немногих столбцов \\ \hline
    q73 & $\approx -41\,\%$ & подсчёт покупателей по числу покупок с фильтрами --- предикат отсекает данные до агрегации \\ \hline
    q88 & $\approx -35\,\%$ & восемь подсчётов с разными временными окнами --- восьмикратное сканирование одного столбца, выигрыш от column pruning умножается \\ \hline
    \multicolumn{3}{|c|}{\textit{запросы, на которых быстрее Parquet}} \\ \hline
    q1 & медленнее & «покупатели, вернувшие товары»: коррелированный подзапрос, агрегация возвратов --- результат во многом определяется стороной соединения и агрегации в JVM \\ \hline
    q2, q28 & медленнее & многократные объединения подзапросов с агрегацией без селективного предиката \\ \hline
    q83, q99 & медленнее & сводки возвратов и заказов с группировкой; преобладает агрегация над декодированными данными \\ \hline
    \end{tabularx}
\end{table}

\textbf{Анализ.} Условия сценария максимально благоприятны для Vortex: таблицы
только для добавления (нет слияния версий), крупные файлы (накладные расходы границы
JNI амортизируются на больших пакетах), реальные данные с избыточностью (кодеки
FastLanes, ALP, FSST дают высокую степень сжатия), запросы с селективными
предикатами и узкими проекциями (нативное проталкивание предикатов сокращает объём
данных, пересекающих границу JVM/native). Пиковое ускорение в \SI{80}{\percent} на
q51 и линейный рост выигрыша с числом сканирований немногих столбцов (q88 --- восемь
сканирований одного столбца) прямо подтверждают первую гипотезу работы.

Запросы, на которых быстрее Parquet, объединяет общая черта --- преобладание
агрегации без селективного предиката над относительно широкими наборами столбцов:
основная работа выполняется не в декодере, а в движке агрегации, который для Parquet
полностью находится в JVM, а для Vortex отделён границей JNI. Это согласуется с
третьей гипотезой и с обсуждавшимся в подразделе~\ref{subsec:methodology}
архитектурным ограничением --- отсутствием проталкивания агрегатов в интерфейсе
форматов Apache Paimon.

Измерить пропускную способность записи на этом наборе не удалось: загрузка крупных
объёмов упиралась в пропускную способность объектного хранилища, и сопоставление
форматов по записи было бы некорректным. Этот аспект покрывается второй серией
экспериментов, где запись выполняется потоком с ограничением скорости.
